\section{Method}
\label{sec:method}

\subsection{Setup}
All runs use one laptop: an Intel Core i7-1355U (2 performance and 8 efficiency cores, 12
hardware threads; AVX2, FMA, F16C and AVX-VNNI, no AVX-512) with 12\,GB of RAM, running Windows
11 (build 26300). The software is the unmodified prebuilt CPU package of llama.cpp release b11696
(\texttt{llama-b11696-bin-win-cpu-x64.zip}, built with Clang 20.1.8). For each ISA tier I made a
folder containing the release with exactly one \texttt{ggml-cpu-*.dll}; the start-up log
confirms which tier loaded. Six tiers load on this CPU: \texttt{x64}, \texttt{sse42},
\texttt{sandybridge} (AVX), \texttt{ivybridge} (AVX, F16C), \texttt{haswell} (AVX2, FMA, F16C)
and \texttt{alderlake} (\texttt{haswell} plus AVX-VNNI). The AVX-512 tiers refuse to load.
\texttt{alderlake} is the tier the loader picks on this machine when all libraries are present.

\paragraph{Models.}
Six GGUF files cover four quantization types and three architectures
(Table~\ref{tab:models}). Q8\_0 and Q4\_K\_M were produced from the F16 file with the release's
\texttt{llama-quantize}. As a control, quantizing on the \texttt{sse42} tier, on the
\texttt{alderlake} tier, and with 1 or 8 threads gave byte-identical files (SHA-256 below), so the
weights themselves do not depend on the machine that made them.

\begin{table}[t]
\centering\small
\caption{Model files. Hashes are the first 16 hex digits of SHA-256. $^{*}$Qwen3.5 mixes
attention layers with linear-recurrence layers.}
\label{tab:models}
\begin{tabular}{llll}
\toprule
label & file & weight types & SHA-256 prefix\\
\midrule
q06-f16   & Qwen3-0.6B F16 & F16 & \texttt{b166332e0224681e}\\
q06-q80   & Qwen3-0.6B Q8\_0 & Q8\_0 & \texttt{f34d7b4d1035d121}\\
q06-q4km  & Qwen3-0.6B Q4\_K\_M & Q4\_K, Q6\_K & \texttt{683e1b8a8ed85096}\\
q06-q40   & Qwen3-0.6B Q4\_0 & Q4\_0 & \texttt{34a119e136d15e41}\\
g3-1b-q40 & Gemma-3-1B-it QAT Q4\_0 & Q4\_0 & \texttt{6364f1139062ff56}\\
q35-08-q40 & Qwen3.5-0.8B Q4\_0 & Q4\_0$^{*}$ & \texttt{f9a8630a4a96076b}\\
\bottomrule
\end{tabular}
\end{table}

\paragraph{Configurations.}
The reference configuration $R$ is: \texttt{alderlake} tier, 8 threads for both generation and
batch processing, logical batch 2048, micro-batch (\texttt{-ub}) 512, flash attention off,
repacking on, one server slot, one client request at a time, prompt cache off. Every other
configuration changes exactly one factor (Table~\ref{tab:configs}). \texttt{server4} is the
setting a parallel evaluation client typically uses against \texttt{llama-server}: four slots,
four concurrent requests and the server's default prompt cache. I ran it twice with identical
commands (\texttt{server4-a}, \texttt{server4-b}) to measure run-to-run variation, and ran $R$
twice for the same purpose.

The two kinds of measurement below cost very different amounts of CPU time. A teacher-forced
run takes about a minute; a benchmark run takes 40 to 55 minutes. So the teacher-forced grid
covers every setting for every model, and the benchmark runs cover a subset chosen in advance
for the primary model (Qwen3-0.6B Q4\_K\_M) and only $R$, one tier change and \texttt{server4}
for the others. I dropped benchmark runs whose teacher-forced result already showed bitwise
agreement with $R$ (\texttt{haswell}, thread counts) and, as the CPU budget ran short, runs that
the grid showed to be numerically equivalent to one already queued. The full list of runs and
the reason for each omission are in \texttt{harness/jobs.py}.

\begin{table}[t]
\centering\small
\caption{Execution-equivalent configurations. Each row differs from the reference in one
factor. The last two columns give the models it was run for in the teacher-forced grid and in
the benchmark runs.}
\label{tab:configs}
\begin{tabular}{llll}
\toprule
name & change from the reference & KL grid & benchmark runs\\
\midrule
ref-rerun & none (second run) & all & Q4\_K\_M\\
isa-haswell, -sse42 & CPU backend tier & all & none\\
isa-ivybridge & tier: AVX, F16C, no FMA & all & F16\\
isa-sandybridge & tier: AVX, no F16C, no FMA & all & Q4\_K\_M, Qwen3 Q4\_0\\
t8, t12 & 8 or 12 threads (KL reference: 4) & all & none\\
ub1, ub8 & micro-batch 1 or 8 & all & none\\
ub7, ub64 & micro-batch 7 or 64 & all & Q4\_K\_M\\
fa-on & flash attention on & all & Q4\_K\_M\\
repack-off & weight repacking off & all & none\\
cache-on & server prompt cache on & none & Q4\_K\_M\\
server4-a & 4 slots, 4 clients, prompt cache on & none & all six\\
server4-b & second run of server4-a & none & Q4\_K\_M\\
\bottomrule
\end{tabular}
\end{table}

\subsection{Evaluation tasks}
I reuse a frozen evaluation set built earlier for an unrelated quantization study with fixed
seeds: zero-shot MMLU \citep{hendrycks2021}, 4-shot chain-of-thought GSM8K \citep{cobbe2021} and
5-shot TriviaQA \citep{joshi2017}. Each configuration runs the first 600 MMLU, 300 TriviaQA and
80 GSM8K items of that file (980 items). Prompts are raw completions with no chat template.
Decoding is greedy (temperature 0, top-$k$ 1). MMLU is scored by reading the top-20 next-token
log-probabilities and taking the highest-scoring answer letter; its \emph{margin} is the
log-probability gap between the best and second-best letter. TriviaQA generates up to 24 tokens
and is scored by normalized alias match; GSM8K generates up to 256 tokens and is scored by the
last number after \texttt{\#\#\#\#}.

\subsection{Measurements}
For every item the harness records the generated token ids, the log-probability of each chosen
token and the first step's top-20 list, exactly as the server prints them (shortest
round-trip decimal of a float32). Against the reference run of the same model I report:
\begin{itemize}
\item \textbf{identical}: token ids, every chosen-token log-probability and the first-step
top-20 list are equal. Equality of these float strings implies equality of the underlying
float32 values.
\item \textbf{same output}: the generated token ids are equal, so the visible answer is equal.
\item \textbf{flip}: the item is scored correct in one run and wrong in the other. I report flips
in both directions, the accuracy change, and an exact McNemar test on the discordant pairs.
\item For generative tasks, the \textbf{first divergence}: the index of the first generated
token that differs.
\end{itemize}

Separately, a teacher-forced grid uses \texttt{llama-perplexity} with \texttt{--kl-divergence}
on 12 chunks of 512 tokens of WikiText; the tool scores the second half of each chunk, so each
run covers 3,072 positions. Its reference is \texttt{alderlake} with 4 threads (so that it could
share the CPU with the benchmark queue), otherwise as $R$. For each model it saves the reference
logits once and then reports, for each configuration, the mean and maximum KL divergence from
the reference and the fraction of positions whose top-1 token agrees. The tool stores the
reference log-probabilities in a compressed 16-bit form, so two bitwise-identical runs do not
read exactly zero: they show $|\text{mean KLD}|\le10^{-6}$ and a maximum around $6\times10^{-5}$.
I call this the storage floor. A configuration at the floor with 100\% top-1 agreement is
reported as indistinguishable from the reference; exact bit equality is established only by the
benchmark runs.

\subsection{Pre-registration}
Before reading any result I wrote down predictions derived from the source code
(\texttt{PREDICTIONS.md}, committed at the start of the run; summarized in
Section~\ref{sec:mechanism}). The design is one-factor-at-a-time; I did not search for
configurations that maximize differences.
